% Budget: 1.5 pages. Claim: Tier 3. Every number here comes from code/e4_scan.py via generated/e4_*.tex.
\section*{E4. Long-horizon scan (Tier 3 claimed)}

\textbf{What is scanned.} All 38 directed backbone links (each direction separately: the photon test uses the sender
at emission and the receiver at arrival), for $t_e \in [0, \EfourYears]$ Julian years ($1{,}753{,}200$ h). $m(t_e)$
is the Sun's distance to the segment from the sender at $t_e$ to the receiver at its light-time arrival; the link is
\emph{closed} iff $m < 0.10$ AU. Maintenance is one-time ($h<264$, handled in S1); the scan covers geometry only.

\textbf{Step size and the shortest closure it could miss.} Base step $\Delta t = \EfourStep$ h. Each segment point
is a fixed convex combination of the endpoints, so if both move at most $\varepsilon$, $m$ changes at most
$\varepsilon$: $|dm/dt_e| \le L$, the larger endpoint speed (perihelion, times $(1+v_{tx}/c)/(1-v_{rx}/c)$ for the
stretch of $t_a$); the largest is $L=\EfourLmax$ AU/h (Mercury). On any $[a,b]$ with end values $m_a, m_b$,
\[ \tfrac12\bigl(m_a+m_b-L(b-a)\bigr) \;\le\; m(t) \;\le\; \tfrac12\bigl(m_a+m_b+L(b-a)\bigr). \]
A closure between two grid points must have grid values within $L\Delta t/2 = \EfourMargin$ AU of the threshold,
last under 1 h, and be shallower than $\EfourMargin$ AU. Every interval whose bound straddles 0.10 AU is bisected until
certified open, certified closed, or shorter than \EfourTolMs\ ms (\EfourEvals\ evaluations), so the scan misses no
closure or reopening longer than \EfourTolMs\ ms; \EfourSlivers\ sub-millisecond intervals, all beside a located
boundary, stay undetermined. The shortest closure found lasts \EfourShortestH\ h (\EfourShortestLink): a step of a
few hours would have been safe, and the certificate makes that a proof within the model.

\textbf{Analytical bounds.} Relays A and B ride one circle of radius $\sqrt8$ AU, $90^\circ$ apart, so their chord
stays $\sqrt8\cos45^\circ = 2$ AU from the Sun (flight motion shifts it under 0.01 AU): A$\leftrightarrow$B never
closes. A gateway's relay link closes only when that relay is almost exactly behind the Sun, which cannot happen to
both at once; the scan's longest simultaneous closure of both is \EfourBothClosedMaxH\ h. \textbf{Every settlement
can always reach every other in at most three links} (open relay plus A$\leftrightarrow$B): geometric connectivity
never fails in 200 years (\EfourNeverCut).

\begin{table}[h]\centering
\gen{e4_links.tex}
\caption{Closures per settlement link over \EfourYears\ Julian years (worse of the two directions). Open \% is the
long-run fraction of emission times at which the link can launch; d = closure duration in days.}
\end{table}

\textbf{Refined boundary (Tier 2).} The first closure on an Earth link is \EfourBndLink. The 1 h grid brackets it in
$[h\,\EfourBndLo, h\,\EfourBndHi]$ with $m = \EfourBndMlo$ and $\EfourBndMhi$ AU. After \EfourBndIter\ bisection steps
the boundary is $t_e = h\,\EfourBndT$ to within \EfourBndWidthMs\ ms (tolerance \EfourTolMs\ ms, the same as the
light-time solution). The closure lasts \EfourBndDays\ days. \emph{Packets in flight:} the test applies at launch,
so the last packet emitted before the boundary flies its full \EfourBndFlight\ min and arrives. Its hop receipt goes
back on \EfourBndRevLink, which closes \EfourBndOffsetMin\ min later (h \EfourBndRevStart), so it still launches
(receipt blocked: \EfourBndReceiptBlocked). The risky case is the \emph{other} direction: a packet launched on
\EfourBndRevLink\ in those \EfourBndOffsetMin\ min arrives, but its receipt needs \EfourBndLink, already closed for
\EfourBndDays\ days. The sender's hop timer ($R_h \approx \EfourBndRh$ min) expires, its retry cannot launch, and it
learns of delivery only after reopening; the relay keeps the packet ID, so a late retry is a recognised duplicate.
This wait fits the 30-day packet lifetime, but at Mars, Jupiter, Saturn and Ceres the longest closures exceed it.
Hence R16: never launch on a link whose reverse closes within one $R_h$, and re-route any send a known closure
would hold more than 24 h.

\begin{table}[h]\centering
\gen{e4_settlements.tex}
\caption{Fastest-open-route one-way delay (empty queues) and worst 24 h availability over \EfourYears\ Julian years
(hourly starts, all four routes of $\le 3$ links). \emph{Naive pin}: fastest route at window start, held 24 h.
\emph{Foresight pin}: route with the best availability over the next 24 h, computable from published geometry.}
\end{table}

\textbf{Poorest service (Tier 3).} By delay: \EfourWorstDelayPair, \EfourWorstDelay\ h one-way (window from hour
\EfourWorstDelayAt). By availability: a naively pinned route can lose the whole window (\EfourWorstNaivePair,
\EfourWorstNaive\% from hour \EfourWorstNaiveAt; every settlement has such windows), because the fastest route can
close minutes later, for \EfourShortestH\ h up to \EfourLongestD\ days (\EfourLongestLink). With foresight pinning,
every ordered pair keeps \EfourWorstFore\% availability in every 24 h window of the 200 years; hence R15.
\textbf{Limits.} These are statements about the common Kepler model, not a forever ephemeris. Random loss (E2) and
incidents (S2) are excluded.
